% GENERATED FILE. Edit canonical lessons, then run npm run generate:books.
% canonical-source-hash: fnv1a64:cedf22e384f3cc17
% canonical-lessons: TE-C520-kudutapadu, TE-C520-balapadu, TE-C520-naligipovu, TE-C520-cidu, TE-C520-punjuko

\chapter{Wearing Down and Bouncing Back}
\label{ch:te-wearing-down-and-bouncing-back}
\hlchaptermodality{520}

\begin{chapteropening}
\textbf{By the end of this chapter:} \emph{I can say that someone got better, grew strong or was worn out.}

Complete the last lesson of chapter 520: I can say that someone got better, grew strong or was worn out.
\end{chapteropening}

\section[\texorpdfstring{\te{కుదుటపడు} (kuduṭapaḍu)}{కుదుటపడు (kuduṭapaḍu)}]{\te{కుదుటపడు} (kuduṭapaḍu) — to steady; to settle down}

\label{lesson:TE-C520-kudutapadu}



\begin{quote}
Before the new one: say the Telugu for a bandage, then the Telugu for to visit the sick.
\end{quote}

\subsection*{You'll want to know: \te{కుదుటపడు}}
\textbf{\te{కుదుటపడు}} — \emph{kuduṭapaḍu} — \textquotedblleft{}to steady; to settle down\textquotedblright{}.

\begin{quote}
\te{మందులతో} \te{అతని} \te{ఆరోగ్యం} \te{కుదుటపడింది}.
\end{quote}

\begin{quote}
\emph{mandulatō atani ārōgyaṁ kuduṭapaḍindi.}
\end{quote}

\begin{quote}
With medicines his health settled down.
\end{quote}

Past (I …): \te{కుదుటపడ్డాను}, \emph{kuduṭapaḍḍānu}.

Health, a situation or a mind settles after trouble.

\begin{cousinweb}
\textbf{\te{కుదుటపడు}} (\emph{kuduṭapaḍu}) is native Dravidian: DEDR 1709 lists Telugu \emph{kuduru}, \textquotedblleft{}to become settled, to become calm, to be cured\textquotedblright{}, and \emph{kuduṭabaḍu} itself, with Tamil \emph{kutir}, \textquotedblleft{}to be settled\textquotedblright{}.
\end{cousinweb}

\begin{grammarlens}[title={What you've built}]
One more word for describing what you lived and felt.
\end{grammarlens}

\subsection*{Your turn}
\begin{itemize}
\raggedright
  \item \emph{Say it:} \emph{kuduṭapaḍu}
  \item \emph{Say it:} \emph{kuduṭapaḍu}, once more
  \item \emph{Say it:} \emph{mandulatō atani ārōgyaṁ kuduṭapaḍindi.} — With medicines his health settled down.
  \item \emph{Recall:} say the Telugu for a bandage, then the Telugu for to visit the sick, then say \emph{kuduṭapaḍu} again
\end{itemize}

\begin{tcolorbox}[breakable,colback=teal!4,colframe=teal!35!black,title={Before you move on}]
What does \te{కుదుటపడు} mean? (\textquotedblleft{}To steady; to settle down\textquotedblright{}.) Say it once more.
\end{tcolorbox}
\section[\texorpdfstring{\te{బలపడు} (balapaḍu)}{బలపడు (balapaḍu)}]{\te{బలపడు} (balapaḍu) — to grow strong}

\label{lesson:TE-C520-balapadu}



\begin{quote}
Before the new one: say the Telugu for to visit the sick, then the Telugu for to steady.
\end{quote}

\subsection*{You'll want to know: \te{బలపడు}}
\textbf{\te{బలపడు}} — \emph{balapaḍu} — \textquotedblleft{}to grow strong\textquotedblright{}.

\begin{quote}
\te{వ్యాయామంతో} \te{నా} \te{కాళ్ళు} \te{బలపడ్డాయి}.
\end{quote}

\begin{quote}
\emph{vyāyāmantō nā kāḷḷu balapaḍḍāyi.}
\end{quote}

\begin{quote}
With exercise my legs grew strong.
\end{quote}

Past (I …): \te{బలపడ్డాను}, \emph{balapaḍḍānu}.

\te{బలం}, strength, + \te{పడు}. Also of friendships: \te{స్నేహం} \te{బలపడింది}.

\begin{cousinweb}
\textbf{\te{బలపడు}} (\emph{balapaḍu}) is built from \te{బలం}, from Sanskrit \emph{bala}, \textquotedblleft{}power, strength\textquotedblright{} (CDIAL 9161) and the native verb \emph{paḍu}, \textquotedblleft{}to feel, to suffer\textquotedblright{} (DEDR 3853).
\end{cousinweb}

\begin{grammarlens}[title={What you've built}]
One more word for describing what you lived and felt.
\end{grammarlens}

\subsection*{Your turn}
\begin{itemize}
\raggedright
  \item \emph{Say it:} \emph{balapaḍu}
  \item \emph{Say it:} \emph{balapaḍu}, once more
  \item \emph{Say it:} \emph{vyāyāmantō nā kāḷḷu balapaḍḍāyi.} — With exercise my legs grew strong.
  \item \emph{Recall:} say the Telugu for to visit the sick, then the Telugu for to steady, then say \emph{balapaḍu} again
\end{itemize}

\begin{tcolorbox}[breakable,colback=teal!4,colframe=teal!35!black,title={Before you move on}]
What does \te{బలపడు} mean? (\textquotedblleft{}To grow strong\textquotedblright{}.) Say it once more.
\end{tcolorbox}
\section[\texorpdfstring{\te{నలిగిపోవు} (naligipōvu)}{నలిగిపోవు (naligipōvu)}]{\te{నలిగిపోవు} (naligipōvu) — to be worn out; to get crumpled}

\label{lesson:TE-C520-naligipovu}



\begin{quote}
Before the new one: say the Telugu for to steady, then the Telugu for to grow strong.
\end{quote}

\subsection*{You'll want to know: \te{నలిగిపోవు}}
\textbf{\te{నలిగిపోవు}} — \emph{naligipōvu} — \textquotedblleft{}to be worn out; to get crumpled\textquotedblright{}.

\begin{quote}
\te{రోజంతా} \te{పనిచేసి} \te{అమ్మ} \te{నలిగిపోయింది}.
\end{quote}

\begin{quote}
\emph{rōjantā panicēsi amma naligipōyindi.}
\end{quote}

\begin{quote}
After working all day, mum was worn out.
\end{quote}

Past (I …): \te{నలిగిపోయాను}, \emph{naligipōyānu}.

Of cloth, to crumple; of people, to be worn down by work or worry.

\begin{cousinweb}
\textbf{\te{నలిగిపోవు}} (\emph{naligipōvu}) is native Dravidian: DEDR 3611 lists Telugu \emph{nalāgu}, \emph{naliyu}, \textquotedblleft{}to be crushed, bruised, crumpled\textquotedblright{}, and the verb ends in \emph{pōvu}, \textquotedblleft{}to go\textquotedblright{} (DEDR 4572).
\end{cousinweb}

\begin{grammarlens}[title={What you've built}]
One more word for describing what you lived and felt.
\end{grammarlens}

\subsection*{Your turn}
\begin{itemize}
\raggedright
  \item \emph{Say it:} \emph{naligipōvu}
  \item \emph{Say it:} \emph{naligipōvu}, once more
  \item \emph{Say it:} \emph{rōjantā panicēsi amma naligipōyindi.} — After working all day, mum was worn out.
  \item \emph{Recall:} say the Telugu for to steady, then the Telugu for to grow strong, then say \emph{naligipōvu} again
\end{itemize}

\begin{tcolorbox}[breakable,colback=teal!4,colframe=teal!35!black,title={Before you move on}]
What does \te{నలిగిపోవు} mean? (\textquotedblleft{}To be worn out; to get crumpled\textquotedblright{}.) Say it once more.
\end{tcolorbox}
\section[\texorpdfstring{\te{చీదు} (cīdu)}{చీదు (cīdu)}]{\te{చీదు} (cīdu) — to blow one's nose}

\label{lesson:TE-C520-cidu}



\begin{quote}
Before the new one: say the Telugu for to grow strong, then the Telugu for to be worn out.
\end{quote}

\subsection*{You'll want to know: \te{చీదు}}
\textbf{\te{చీదు}} — \emph{cīdu} — \textquotedblleft{}to blow one's nose\textquotedblright{}.

\begin{quote}
\te{జలుబుతో} \te{అతను} \te{పదే} \te{పదే} \te{ముక్కు} \te{చీదాడు}.
\end{quote}

\begin{quote}
\emph{jalubutō atanu padē padē mukku cīdāḍu.}
\end{quote}

\begin{quote}
With his cold he blew his nose again and again.
\end{quote}

Past (I …): \te{చీదాను}, \emph{cīdānu}.

With \te{ముక్కు}, nose: \te{ముక్కు} \te{చీదు}.

\begin{cousinweb}
\textbf{\te{చీదు}} (\emph{cīdu}) is native Dravidian: DEDR 2618 lists the Telugu verb with Tamil \emph{cīntu}, \textquotedblleft{}to blow, as the nose\textquotedblright{}.
\end{cousinweb}

\begin{grammarlens}[title={What you've built}]
One more word for describing what you lived and felt.
\end{grammarlens}

\subsection*{Your turn}
\begin{itemize}
\raggedright
  \item \emph{Say it:} \emph{cīdu}
  \item \emph{Say it:} \emph{cīdu}, once more
  \item \emph{Say it:} \emph{jalubutō atanu padē padē mukku cīdāḍu.} — With his cold he blew his nose again and again.
  \item \emph{Recall:} say the Telugu for to grow strong, then the Telugu for to be worn out, then say \emph{cīdu} again
\end{itemize}

\begin{tcolorbox}[breakable,colback=teal!4,colframe=teal!35!black,title={Before you move on}]
What does \te{చీదు} mean? (\textquotedblleft{}To blow one's nose\textquotedblright{}.) Say it once more.
\end{tcolorbox}
\section[\texorpdfstring{\te{పుంజుకో} (puñjukō)}{పుంజుకో (puñjukō)}]{\te{పుంజుకో} (puñjukō) — to pick up; to regain strength}

\label{lesson:TE-C520-punjuko}



\begin{quote}
Before the new one: say the Telugu for to be worn out, then the Telugu for to blow one's nose.
\end{quote}

\subsection*{You'll want to know: \te{పుంజుకో}}
\textbf{\te{పుంజుకో}} — \emph{puñjukō} — \textquotedblleft{}to pick up; to regain strength\textquotedblright{}.

\begin{quote}
\te{జ్వరం} \te{తగ్గాక} \te{పాప} \te{మళ్ళీ} \te{పుంజుకుంది}.
\end{quote}

\begin{quote}
\emph{jvaraṁ taggāka pāpa maḷḷī puñjukundi.}
\end{quote}

\begin{quote}
After the fever went down, the little girl picked up again.
\end{quote}

Past (I …): \te{పుంజుకున్నాను}, \emph{puñjukunnānu}.

Said of a person getting their strength back, and of a team or a business that rallies.

\begin{cousinweb}
\textbf{\te{పుంజుకో}} (\emph{puñjukō}) is a Telugu verb ending in the reflexive -\te{కో}; the origin of its stem is not settled.
\end{cousinweb}

\begin{grammarlens}[title={What you've built}]
One more word for describing what you lived and felt.
\end{grammarlens}

\subsection*{Your turn}
\begin{itemize}
\raggedright
  \item \emph{Say it:} \emph{puñjukō}
  \item \emph{Say it:} \emph{puñjukō}, once more
  \item \emph{Say it:} \emph{jvaraṁ taggāka pāpa maḷḷī puñjukundi.} — After the fever went down, the little girl picked up again.
  \item \emph{Recall:} say the Telugu for to be worn out, then the Telugu for to blow one's nose, then say \emph{puñjukō} again
\end{itemize}

\begin{tcolorbox}[breakable,colback=teal!4,colframe=teal!35!black,title={Before you move on}]
What does \te{పుంజుకో} mean? (\textquotedblleft{}To pick up; to regain strength\textquotedblright{}.) Say it once more.
\end{tcolorbox}
